\chapter{Mathematical Preliminaries: Function Spaces}

\section{Hölder Spaces}

\begin{definition}
Let $U\subset \mathbb{R}^n$ be open, let $0<\gamma\le 1$, and let
$
u:U\to \mathbb{R}.
$
Suppose $u$ is bounded and continuous. We define
$$
\|u\|_{C(\overline{U})}
:=
\sup_{x\in U}|u(x)|.
$$
\end{definition}

\begin{definition}[Hölder continuity]
We say that $u$ is \emph{Hölder continuous with exponent $\gamma$} if there exists a constant $C>0$ such that
$$
|u(x)-u(y)|
\le
C|x-y|^\gamma
$$
for all $x,y\in U$.
\end{definition}
\begin{remark}\
    \begin{itemize}
        \item Lipschitz continuity is exactly Hölder continuity with exponent is 1.
    \end{itemize}
\end{remark}

\begin{definition}[Hölder seminorm and Hölder norm]
The $\gamma$-th Hölder seminorm of $u$ is defined by
$$
[u]_{C^{0,\gamma}(\overline{U})}
:=
\sup_{\substack{x,y\in U\\ x\neq y}}
\left\{
\frac{|u(x)-u(y)|}{|x-y|^\gamma}
\right\}.
$$

The $\gamma$-th Hölder norm is defined by
$$
\|u\|_{C^{0,\gamma}(\overline{U})}
:=
\|u\|_{C(\overline{U})}
+
[u]_{C^{0,\gamma}(\overline{U})}.
$$
\end{definition}

\begin{definition}[Hölder spaces]
Let $k\in \mathbb{N}$ and $0<\gamma\le 1$. The Hölder space is defined as
$$
C^{k,\gamma}(\overline{U})
:=
\left\{
u\in C^k(\overline{U}):
\|u\|_{C^{k,\gamma}(\overline{U})}<\infty
\right\}.
$$
The norm is
$$
\|u\|_{C^{k,\gamma}(\overline{U})}
:=
\sum_{|\alpha|\le k}
\|D^\alpha u\|_{C(\overline{U})}
+
\sum_{|\alpha|=k}
[D^\alpha u]_{C^{0,\gamma}(\overline{U})}.
$$
\end{definition}

\begin{theorem}
The Hölder space
$
C^{k,\gamma}(\overline{U})
$
is a Banach space.
\end{theorem}










\section{$L^p$ Spaces}

\begin{definition}
Let $U\subset \mathbb{R}^n$ be measurable and let
$
1\le p\le \infty.
$
We define
$$
L^p(U)
:=
\left\{
f:U\to \mathbb{R}:
f \text{ is measurable and } \|f\|_{L^p(U)}<\infty
\right\}.
$$
For $1\le p<\infty$, the $L^p$ norm is
$$
\|f\|_{L^p(U)}
:=
\left(
\int_U |f(x)|^p\,dx
\right)^{1/p}.
$$
For $p=\infty$, the $L^\infty$ norm is
$$
\|f\|_{L^\infty(U)}
:=
\operatorname*{ess\,sup}_{x\in U}|f(x)|.
$$
\end{definition}

\begin{theorem}[Minkowski inequality]
For $1\le p\le \infty$,
$$
\|f+g\|_{L^p(U)}
\le
\|f\|_{L^p(U)}+\|g\|_{L^p(U)}.
$$
\end{theorem}

\begin{theorem}[Hölder inequality]
Let
$$
1\le p,q\le \infty,
\qquad
\frac{1}{p}+\frac{1}{q}=1.
$$
Then
$$
\|fg\|_{L^1(U)}
\le
\|f\|_{L^p(U)}\|g\|_{L^q(U)}.
$$
\end{theorem}

\begin{theorem}
The spaces $L^p(U)$ are Banach spaces.
\end{theorem}









\section{Weak Derivatives and Sobolev Spaces}

\begin{definition}[Weak derivative]
Let
$
u,v\in L^1_{\mathrm{loc}}(U),
$
and let $\alpha$ be a multi-index.
We say that $v$ is the $\alpha$-th weak partial derivative of $u$ if
$$
\int_U u\,D^\alpha \phi\,dx
=
(-1)^{|\alpha|}
\int_U v\,\phi\,dx
$$
for all test functions
$
\phi\in C_c^\infty(U).
$
We denote this by
$$
D^\alpha u=v.
$$
\end{definition}

\begin{example}
See Example 1 and Example 2.
\end{example}

\begin{remark}
    \textbf{What is the relation between weak derivative and the distributional derivative?}
    \begin{itemize}
        \item 
        Let $\Omega \subset \mathbb{R}^n$ be open and let
        $u \in L^1_{\mathrm{loc}}(\Omega)$.
        
        A locally integrable function $u$ defines a distribution $T_u$ by
        \[
        T_u(\varphi)
        =
        \int_\Omega u(x)\varphi(x)\,dx,
        \qquad
        \varphi \in C_c^\infty(\Omega).
        \]
        
        The \emph{distributional derivative} of $u$ in the $x_i$-direction is
        the distribution $\partial_i T_u$ defined by
        \[
        \langle \partial_i T_u,\varphi\rangle
        =
        -\langle T_u,\partial_i\varphi\rangle
        =
        -\int_\Omega u(x)\partial_i\varphi(x)\,dx.
        \]

        
        Thus the distributional derivative is a functional on test functions.
        
        A function $v \in L^1_{\mathrm{loc}}(\Omega)$ is called the
        \emph{weak derivative} of $u$ in the $x_i$-direction if
        \[
        \int_\Omega v(x)\varphi(x)\,dx
        =
        -\int_\Omega u(x)\partial_i\varphi(x)\,dx
        \]
        for every $\varphi \in C_c^\infty(\Omega)$.
        In terms of distributions, this condition is
        \[
        T_v = \partial_i T_u.
        \]
        
        Therefore,
        \begin{itemize}
            \item $
                v \text{ is the weak derivative of } u
                \iff
                T_v \text{ is the distributional derivative of } T_u.
                $
            \item $
                \text{weak derivative}
                =
                \text{function representative of the distributional derivative}.
                $
        \end{itemize}
        
        \item
        Not every distributional derivative is represented by a function.
        For example, let
        \[
        H(x)
        =
        \begin{cases}
        0, & x<0,\\
        1, & x>0.
        \end{cases}
        \]
        Then
        \[
        H' = \delta_0
        \]
        in the sense of distributions, where
        \[
        \delta_0(\varphi)=\varphi(0).
        \]
        Since $\delta_0$ is not given by integration against a locally integrable
        function, $H$ has a distributional derivative but no weak derivative as an
        ordinary function.
        
        In summary, Every weak derivative gives a distributional derivative,
        but not every distributional derivative is a weak derivative.

    \end{itemize}
\end{remark}

\begin{proposition}[Uniqueness of weak derivatives]
A weak derivative, if it exists, is uniquely defined up to a set of measure zero.
\end{proposition}

\begin{definition}[Sobolev spaces]
Let
$
U\subset \mathbb{R}^n
$
be open, let
$
1\le p<\infty,
$
$
k\in \mathbb{N}\cup\{0\}.
$
The Sobolev space
$
W^{k,p}(U)
$
consists of all locally summable functions
$
u:U\to \mathbb{R}
$
such that for each multi-index $\alpha$ with
$
|\alpha|\le k,
$
the weak derivative $D^\alpha u$ exists in the weak sense and belongs to $L^p(U)$.
\end{definition}

\begin{remark}
Equivalently,
$$
W^{k,p}(U)
=
\left\{
u\in L^p(U):
D^\alpha u\in L^p(U)
\text{ for all } |\alpha|\le k
\right\},
$$
where $D^\alpha u$ is the distributional derivative of $u$.
\end{remark}

\begin{theorem}[Properties of weak derivatives]
Assume
$
u,v\in W^{k,p}(U),
\text{ with }
|\alpha|\le k.
$
Then:
\begin{enumerate}
    \item
    $
    D^\alpha u\in W^{k-|\alpha|,p}(U).
    $
    Moreover,
    $$
    D^\beta(D^\alpha u)
    =
    D^\alpha(D^\beta u)
    =
    D^{\alpha+\beta}u
    $$
    for all multi-indices $\alpha,\beta$ with
    $
    |\alpha|+|\beta|\le k.
    $

    \item For all $\lambda,\mu\in \mathbb{R}$,
    $$
    \lambda u+\mu v\in W^{k,p}(U),
    $$
    and
    $$
    D^\alpha(\lambda u+\mu v)
    =
    \lambda D^\alpha u+\mu D^\alpha v.
    $$

    \item If $V$ is an open subset of $U$, then
    $
    u\in W^{k,p}(V).
    $

    \item If
    $
    \zeta\in C_c^\infty(U),
    $
    then
    $
    \zeta u\in W^{k,p}(U),
    $
    and
    $$
    D^\alpha(\zeta u)
    =
    \sum_{\beta\le \alpha}
    \binom{\alpha}{\beta}
    D^\beta \zeta\,
    D^{\alpha-\beta}u.
    $$
    This is the Leibniz formula.
\end{enumerate}
\end{theorem}

\begin{remark}\
    Here
    $
    \beta\le \alpha
    $
    means that
    $
    \beta_i\le \alpha_i
    $
    for each
    $
    i=1,\ldots,n,
    $
    and
    $$
    \binom{\alpha}{\beta}
    :=
    \frac{\alpha!}{\beta!(\alpha-\beta)!},
    $$
    where
    $
    \alpha!=(\alpha_1!)\cdots(\alpha_n!).
    $
\end{remark}


\begin{proof}
Part (4) follows by induction on $|\alpha|$.
\end{proof}

\begin{definition}[Sobolev norm]
    For $1\le p\le \infty$, define the Sobolev norm by
    $$
    \|u\|_{W^{k,p}(U)}
    =
    \begin{cases}
    \left(
    \displaystyle\sum_{|\alpha|\le k}
    \int_U |D^\alpha u|^p\,dx
    \right)^{1/p},
    & 1\le p<\infty,\\[1.2em]
    \displaystyle\sum_{|\alpha|\le k}
    \operatorname*{ess\,sup}_{U}|D^\alpha u|,
    & p=\infty.
    \end{cases}
    $$
    \end{definition}
    
    \begin{example}
    For $k=1$ and $1\le p<\infty$,
    $$
    \|u\|_{W^{1,p}(U)}
    =
    \left(
    \int_U |u|^p\,dx
    +
    \sum_{i=1}^n
    \int_U
    \left|
    \partial_i u
    \right|^p
    \,dx
    \right)^{1/p}
    =
    \left(
    \|u\|_{L^p(U)}^p
    +
    \|Du\|_{L^p(U)}^p
    \right)^{1/p}.
    $$
    \end{example}

    \begin{remark}\
        \begin{itemize}
            \item $W^{k,p}(U)$, with $k\ge 1$ and $1\le p\le \infty$, is a Banach space.
            \item $H^k(U) := W^{k,2}(U)$ is actually a Hilbert space.
        \end{itemize}
    \end{remark}


    \section{Approximation of Sobolev Functions}

\begin{definition}[Convergence in Sobolev spaces]
Let
$
\{u_m\}_{m=1}^{\infty}\subset W^{k,p}(U),
$
$
u\in W^{k,p}(U).
$
We say that $u_m$ converges to $u$ in $W^{k,p}(U)$ if
$$
\lim_{m\to\infty}
\|u_m-u\|_{W^{k,p}(U)}
=0.
$$
We denote this by
$
u_m\to u
$
in
$W^{k,p}(U).
$
\end{definition}

\begin{definition}[Local Sobolev space]
The local Sobolev space is defined by
$$
W^{k,p}_{\mathrm{loc}}(U)
:=
\left\{
u:
u\in W^{k,p}(V)
\text{ for every } V\subset\subset U
\right\}.
$$
We write
$
u_m\to u
$
in
$W^{k,p}_{\mathrm{loc}}(U)
$
if
$
u_m\to u
$
in
$W^{k,p}(V)
$
for each
$
V\subset\subset U.
$
\end{definition}

\begin{theorem}[Local approximation by smooth functions]
Assume
$
u\in W^{k,p}(U)
$
for some
$
1\le p<\infty.
$
Set
$
u^\varepsilon:=\eta_\varepsilon*u
$
in
$U_\varepsilon.
$
Then:
\begin{enumerate}
    \item
    $
    u^\varepsilon\in C^\infty(U_\varepsilon)
    \qquad
    \forall \varepsilon>0.
    $

    \item
    $
    u^\varepsilon\to u
    \text{ in } W^{k,p}_{\mathrm{loc}}(U)
    $
    as
    $
    \varepsilon\to 0.
    $
\end{enumerate}
\end{theorem}

\begin{theorem}[Meyers--Serrin theorem]
    Let $U\subset \mathbb{R}^n$ be open and bounded. Then
    $
    C^\infty(U)\cap W^{k,p}(U)
    $
    is dense in
    $
    W^{k,p}(U).
    $
    \end{theorem}
    
    \begin{remark}\
    \begin{itemize}
        \item This is the global approximation by smooth functions.
        \item Equivalently, we can say that there exist functions
        $
        u_m\in C^\infty(U)\cap W^{k,p}(U)
        $
        such that
        $
        u_m\to u
        \text{ in } W^{k,p}(U).
        $
    \end{itemize}
    \end{remark}

\begin{theorem}[Global approximation by functions smooth up to the boundary]
    Assume $U$ is bounded and $\partial U$ is $C^1$. Suppose
    $
    u\in W^{k,p}(U)
    $
    for some
    $
    1\le p<\infty.
    $
    Then there exist functions
    $
    u_m\in C^\infty(\overline{U})
    $
    such that
    $
    u_m\to u
    \text{ in } W^{k,p}(U).
    $
\end{theorem}


\begin{theorem}[Extension theorem]
    Assume $U$ is bounded and $\partial U$ is $C^1$. Select a bounded open set $V$ such that
    $
    U\subset\subset V.
    $
    Then there exists a bounded linear operator
    $$
    E:W^{1,p}(U)\to W^{1,p}(\mathbb{R}^n)
    $$
    such that for each
    $
    u\in W^{1,p}(U),
    $
    we have:
    \begin{enumerate}
        \item
        $
        Eu=u \text{ a.e. in } U;
        $
    
        \item
        $
        \operatorname{supp}(Eu)\subset V;
        $
    
        \item
        $
        \|Eu\|_{W^{1,p}(\mathbb{R}^n)}
        \le
        C\|u\|_{W^{1,p}(U)},
        $
        where the constant $C$ depends only on $p$, $U$, and $V$.
    \end{enumerate}
    \end{theorem}
    
    \begin{definition}
    We call $Eu$ an \emph{extension} of $u$ to $\mathbb{R}^n$.
    \end{definition}









\section{Embeddings Between Sobolev Spaces}
If
$
f\in W^{k,p}(U),
$
does this automatically imply that $f$ belongs to any other spaces?

\begin{theorem}[Gagliardo--Nirenberg--Sobolev inequality]
Assume
$
1\le p<n.
$
Then there exists a constant
$
C=C(n,p)
$
such that, for all
$
u\in C_c^1(\mathbb{R}^n),
$
we have
$$
\|u\|_{L^q(\mathbb{R}^n)}
\le
C\|Du\|_{L^p(\mathbb{R}^n)},
$$
where
$
\frac{1}{q}+\frac{1}{n}=\frac{1}{p}.
$
Equivalently,
$
q=\frac{np}{n-p}.
$
\end{theorem}

\begin{corollary}[Estimate for $W^{1,p}$, $1\leq p <\infty$]
Let
$
U\subset \mathbb{R}^n
$
be bounded and open, and suppose
$
\partial U \text{ is } C^1.
$
Assume
$
1\le p<n
$
and
$
u\in W^{1,p}(U).
$
Then
$
u\in L^q(U),
$
with the estimate
$$
\|u\|_{L^q(U)}
\le
C\|u\|_{W^{1,p}(U)}.
$$
\end{corollary}

\begin{definition}
Let
$
U\subset \mathbb{R}^n
$
be open and bounded. We denote
$$
W_0^{1,p}(U)
:=
\overline{C_c^\infty(U)}^{\,W^{1,p}(U)}.
$$
That is, $W_0^{1,p}(U)$ is the closure of $C_c^\infty(U)$ in the $W^{1,p}(U)$ norm.
\end{definition}

\begin{corollary}[Estimate for $W_0^{1,p}$, $1\leq p <\infty$]
Let
$
U\subset \mathbb{R}^n
$
be bounded and open. Suppose
$
u\in W_0^{1,p}(U)
$
for some
$
1\le p<n.
$
Then we have the estimate
$$
\|u\|_{L^q(U)}
\le
C\|Du\|_{L^p(U)}
$$
for each
$
q\in \left[1,\frac{np}{n-p}\right],
$
where the constant
$
C=C(q,p,n,U).
$
\end{corollary}

\begin{remark}\
    \begin{itemize}
        \item For $n=1$, the restriction $p<n$ is necessary.
        For example,
        $$
        x\mapsto \log\!\left(\log\!\left(1+\frac{1}{|x|}\right)\right)
        $$
        belongs to
        $
        W^{1,n}(B_1),
        $
        but not to
        $
        L^\infty(B_1).
        $

        \item For the case
        $
        n>p,
        $
        we use Morrey's inequality.
    \end{itemize}
\end{remark}